\documentclass{article}

\usepackage{/home/paul/Workspace/build/macros}

\title{Devoir maison}
\author{Paul Catala}

\usepackage[T1]{fontenc}

\usepackage{matlab-prettifier}
\usepackage{listings}
\lstset{
	style = Matlab-editor,
	breaklines = true,
	basicstyle = \footnotesize\ttfamily,
	extendedchars = true,
	frame = single
}

\usepackage{geometry}
\geometry{letterpaper}    
\usepackage{graphicx}

\newtheorem{exercise}{Exercise}

%\title{Devoir en optimisation numérique -- ISC 962}

\newenvironment{solution}[1][\it{Solution}]{\textcolor{blue}{\textbf{#1. } }}{\textcolor{blue}{$\blacksquare$}}


\begin{document}
\noindent Master Ingénierie des Systèmes Complexes \hfill M2 ISC 903.4 \\
Optimisation numérique et sciences des données \hfill 13/01/2025 


\begin{center}
\large{Devoir en pratique de l'optimisation numérique -- M2 ISC 903.4}
\end{center}

\hrulefill


\bigbreak

\paragraph{Livrables}: Il est demandé de livrer sur Arche au plus tard le \textbf{mardi 04 février 2025 à 18h} un compte-rendu au format pdf du TP ci-dessous.

\paragraph{Consignes}: Le code Matlab devra être commenté au fur et à mesure des questions. Le compte-rendu devra être généré par la commande \texttt{publish} de MATLAB:
\begin{lstlisting}
publish("nom_du_script.m", format="pdf")
\end{lstlisting}

Les figures et autres résultats devront être commentés et interprétés. Indiquer les difficultés que vous avez rencontrées, et vos tentatives pour les surmonter; signaler vos sources d'information. Si vous échangez au sein de la promotion pour avancer dans la résolution d'un problème, il conviendra de citer la ou les personnes qui vous auront permis de progresser.

\newpage

\section*{Résolution d'un problème inverse: déconvolution d'un signal triangulaire}

\subsection*{Objectif de l'exercice}
Le problème considéré est celui de la reconstruction d’un signal $x(t)$ (monodimensionnel), supposé inconnu car non accessible à une mesure directe, ayant subi un filtrage gaussien (convolution) inhérent à l’instrument de mesure. Il s’agit d’un filtrage linéaire d’usage courant en traitement d’image. La réponse impulsionnelle d’un filtre gaussien est définie par :
\eql{\label{eq:filtre}
	h(t) = \frac{1}{\sigma\sqrt{2\pi}} \exp\left(-\frac{-(t-m)^2}{2\sigma^2}\right)
}
où $m$ et $\sigma^2$ sont la moyenne et la variance de la gaussienne, respectivement.

Le signal $y$ mesuré en sortie du filtre, est supposé être corrompu par des incertitudes additives. En effet, les mesures en sortie de filtre sont à minima sujettes au bruit de quantification dû au convertisseur analogique-numérique de la chaîne de mesure. Il s’agit d’appréhender le caractère instable de l’opération de déconvolution, lorsque celle-ci est effectuée de façon \guillemotleft naïve \guillemotright, soit par division polynomiale (fonction \texttt{deconv} de Matlab), soit par filtrage inverse (dans le domaine fréquentiel ou par moindres carrés qui est l’équivalent
matriciel). 
%
En présence d’erreurs (de bruit) sur les données, même très faibles, les résultats fournis par ces méthodes d’inversion peuvent, en effet, être inexploitables. Il s’agit là d’un problème inverse mal posé. La réponse en fréquence du système à inverser permet de caractériser le degré de difficulté du problème. L’exercice illustrera la (les) solution(s) obtenues par régularisation linéaire quadratique.


\subsection*{1. Construction du jeu de données}

La fréquence d'échantillonage est fixée à la valeur $f_e = 1 Hz$. Les échantillons du signal d’entrée sont notés $x[n]$, ceux de la réponse impulsionnelle du filtre $h[n]$ et ceux du signal mesuré en sortie $y[n]$. Le signal mesuré en sortie est sujet à des incertitudes additives notées $b[n]$ et le problème direct peut alors être modélisé comme suit :

\eql{\label{eq:conv}
	y[n] = (x \star h)[n] + b[n]
}

\subsubsection*{a. Simulation du signal d'entrée}
Le signal d'entrée, supposé être triangulaire, est défini par:
\eq{
	x(t) = \left\{\begin{array}{ll}
		\frac{t}{L/2} &\forall t\in[0,\frac{L}{2}]\\
		\frac{-t+L}{L/2} &\forall t \in[L/2,L]
		\end{array}\right.
}
\begin{enumerate}
	\setcounter{enumi}{0}
	\item  Simuler le signal à temps discret $x[n]$ correspondant, pour une durée d'observation $L=100s$.

\textbf{Aide:} remplacer $t$ par $nT_e$ dans l'équation définissant $x(t)$. Il est possible d'utiliser une instruction du type $\texttt{x=[x1,x2]}$, où $\texttt{x1}$ et $\texttt{x2}$ représentent deux tronçon susccessifs du signal. De manière générale, il est conseillé d'adapter le script ci-dessous qui permet de maîtriser les différents paramètres de la discrétisation
\begin{lstlisting}
L = 100;	% dure d'observation du signal d'entree (secondes)
fe = 1;		% frequence d'echantillonage (Hz)
Te = 1/fe;	% periode d'echantillonage (secondes)
Nx = L/Te;	% nombre d'echantillons du signal x
nx = 0:Nx-1;	% indice temporel (entier)
tx = nx*Te;	% base de temps (secondes)
\end{lstlisting}
\end{enumerate}

\subsubsection*{b. Simulation du filtre} 
\begin{enumerate}
	\setcounter{enumi}{1}
\item Pour les besoins de la simulation, déterminer la réponse impulsionnelle $h[n]$ du filtre gaussien obtenue par discrétisation de l’équation \eqref{eq:filtre}. Le filtre est de plus supposé causal (\ie, de valeurs nulles aux temps négatifs) et ses paramètres
seront pris égaux respectivement à $m = 15$ et $\sigma = 4$. Enfin, la réponse impulsionnelle est tronquée à un horizon de $N_h = 30$ s.
	Stocker cette réponse impulsionnelle dans un vecteur et la représenter.
\end{enumerate}

\subsubsection*{c. Simulation du problème direct: première approche}
\begin{enumerate}
	\setcounter{enumi}{4}
	\item Comme le filtre gaussien est un filtre \textbf{linéaire invariant}, il est possible de calculer sa sortie comme le produit de convolution discret donné par l'équation \eqref{eq:conv}. Simuler le signal de sortie non bruité $y_{nb}[n]$ à l'aide de la fonction \texttt{conv}.
	
	\item Si Le signal de sortie $y$ est mesuré à l'aide d'un convertisseur analogique-numérique, il est sujet au bruit de quantification. L'opération de quantification peut être simulée par l'instruction \texttt{y=round}$(y_nb*2^{10})/2^{10}$. Simuler le signal quantifié $y$ et représenter celui-ci à l'aide de la fonction \texttt{stairs}.
	
	\item Déterminer et représenter le bruit de quantification $b[n]$.
\end{enumerate}

\subsubsection*{d. Simulation du problème direct par approche matricielle}
Dans cette partie, le signal de sortie non bruité est simulé par l’équation matricielle $y_{nb} = Hx$. L’usage d’un paramétrage soigné est recommandé pour permettre de modifier aisément les conditions de simulation et d’affichage des résultats.
\begin{enumerate}
	\setcounter{enumi}{7}
	\item Déterminer la dimension $N_y$ du vecteur de sortie $y_{nb}$ à partir des dimensions $N_x$ et $N_h$ des vecteurs d’entrée $x$ et de réponse impulsionnelle $h$ afin d’obtenir le même nombre d’échantillons que par la fonction \texttt{conv}.
	\item Construire la matrice de convolution $H$ à l’aide de la fonction \texttt{toeplitz}. Pour cela, il est proposé d’initialiser le premier vecteur ligne et le premier vecteur colonne de $H$ avec des vecteurs nuls de dimension appropriée à l’aide la fonction \texttt{zeros}, puis de remplacer certains de leurs éléments avec les éléments de $h$.
	\item Déterminer la sortie non bruitée par une opération matricielle, puis appliquer la fonction permettant de simuler l’erreur de troncature.
\end{enumerate}

\subsection*{2. Déconvolution}
\subsubsection*{a. Déconvolution « naïve » et constat de l’instabilité}
\begin{enumerate}
	\setcounter{enumi}{10}
	\item Reconstruire le signal d’entrée en utilisant l’opération de déconvolution de Matlab (\texttt{deconv}), en se plaçant, dans un premier temps, dans le cas idéal où l’on a accès à la sortie non bruitée $y_{nb}$.
	\item A présent, reconstruire le signal d’entrée à partir des échantillons mesurés du signal de sortie $y$ (bruité). Le signal reconstruit sera appelé \texttt{xrec}
	\item Tester et commenter l’influence de la « forme » de la réponse impulsionnelle sur la qualité de la reconstruction en modifiant la valeur du paramètre $\sigma$.
	\end{enumerate}
	
\subsubsection*{b. Déconvolution par moindres carrés ($\ell_2$)}
\begin{enumerate}
	\setcounter{enumi}{13}
	\item Calculer le signal déconvolué à l’aide de l’estimateur des moindres carrés.
	\item Tracer le signal déconvolué et calculer l’énergie de l’erreur de reconstruction.
	\item Tracer les valeurs singulières de la matrice $H$. Que vaut le ratio entre la plus grande valeur singulière, et la plus petite valeur singulière non nulle? Ce quotient est le \guillemotleft conditionnement \guillemotright de $H$, et donne une indication quantitative sur la stabilité de la solution des moindres carrés.
	\item Comment le conditionnement évolue-t-il en fonction de la valeur du paramètre $\sigma$ du filtre?
\end{enumerate}

\subsubsection*{c. Déconvolution par régularisation linéaire quadratique ($\ell_2\ell_2$)}
Cette partie concerne la mise en oeuvre de la solution obtenue par minimisation d'un critère d'optimisation composite $J(x) = F_b(x) + \la Fx(x)$ (problème régularisé)

\begin{enumerate}
		\setcounter{enumi}{17}
		\item Construire la matrice carrée $D_1$, de dimension $N_x$, telle que le produit $D_1x$ donne les différences premières : $x[n+1]-x[n]$ (les différences premières, lorsqu'elles sont divisées par le pas $T_e$, correspondent à une approximation de la dérivée première).

\textbf{Aide} :
\begin{lstlisting}
dcol1 = zeros (1,Nx) ;
dcol1 (1:2) = dcol1 (1:2) + [1 -1];
dlig1 = zeros (1,Nx);
dlig1 (1 ,1) = dcol1(1,1);
D1 = toeplitz (dcol1,dlig1);
\end{lstlisting}
\end{enumerate}


\noindent Pour régulariser les moindrs carrés, on propose de minimiser la fonctionnelle suivante : 
\eq{
	J(x) = \frac{1}{2}\norm{y-Hx}_2^2 + \lambda\norm{D_1 x}_2^2
}
	\begin{enumerate}
		\setcounter{enumi}{18}
		\item Calculer la solution régularisée en déterminant le coefficient de régularisation $\lambda$ de manière empirique  (par exemple en testant plusieurs valeurs dans un intervalle conséquent, comme $[10^{-7},10^2]$).
		
		\item Le signal à reconstruire étant ici parfaitement connu, il est possible de déterminer la valeur du coefficient de régularisation $\la$ qui minimise l'énergie de l'erreur de reconstruction. Tracer l'évolution de l'énergie de l'erreur de reconstruction en fonction de
$\lambda$. Ecrire une procédure permettant d'approcher la valeur optimale de $\lambda$, par recherche exhaustive sur un intervalle, à l'aide de la fonction \texttt{min}.
	\end{enumerate}

\end{document}